% Created 2026-09-21 Mon 11:18
% Intended LaTeX compiler: pdflatex
\documentclass[12pt,a4paper]{article}
\usepackage[utf8]{inputenc}
\usepackage[T1]{fontenc}
\usepackage{graphicx}
\usepackage{grffile}
\usepackage{longtable}
\usepackage{wrapfig}
\usepackage{rotating}
\usepackage[normalem]{ulem}
\usepackage{amsmath}
\usepackage{textcomp}
\usepackage{amssymb}
\usepackage{capt-of}
\usepackage{hyperref}
\usepackage[a4paper,margin=1in]{geometry}
\usepackage{graphicx}
\usepackage{booktabs}
\usepackage{amsmath,amssymb}
\usepackage{float}
\usepackage{hyperref}
\usepackage{xcolor}
\usepackage{listings}
\usepackage{enumitem}
\usepackage{fancyhdr}
\usepackage{longtable}
\hypersetup{colorlinks=true,linkcolor=blue,urlcolor=blue}
\pagestyle{fancy}
\fancyhf{}
\lhead{ICS1512}
\rhead{Machine Learning Algorithms Laboratory}
\cfoot{\thepage}
\lstset{language=Python,basicstyle=\ttfamily\small,numbers=left,frame=single,breaklines=true,showstringspaces=false}
\author{Max}
\date{\today}
\title{ICS1512 -- Machine Learning Laboratory}
\hypersetup{
 pdfauthor={Max},
 pdftitle={ICS1512 -- Machine Learning Laboratory},
 pdfkeywords={},
 pdfsubject={},
 pdfcreator={Emacs 27.1 (Org mode 9.3)}, 
 pdflang={English}}
\begin{document}

\maketitle

\section{Cover}
\label{sec:org7cf6d62}
\begin{center}
\begin{tabular}{|p{4cm}|p{11cm}|}
Experiment No. & 7\\
Experiment Title & Bagging, Boosting, and Stacked Ensemble Models\\
Student Name & Emmanuel Maximus Ravichandran\\
Register Number & 3122247001015\\
Faculty & Dr. Poreddy Ajay Kumar Reddy\\
Submission Date & \\
\end{tabular}
\end{center}

\section{Objective}
\label{sec:orge06a1de}
\begin{itemize}
\item To understand ensemble learning strategies: Bagging, Boosting, and Stacking.
\item To implement Bagging and Boosting classifiers.
\item To build a Stacked Ensemble model using multiple base learners.
\item To compare ensemble models in terms of accuracy, stability, and generalization.
\item To analyze the effect of ensemble methods on bias and variance.
\end{itemize}

\section{Problem Statement}
\label{sec:orga80d4f0}
Given the Wisconsin Diagnostic Breast Cancer (WDBC) dataset of 569 samples described by 30
numerical features computed from digitized images of a breast mass, the task is to classify each
sample as \textbf{Malignant} or \textbf{Benign}. Three ensemble strategies -- Bagging, Boosting, and Stacking --
are implemented, tuned via 5-fold cross-validation, and compared to determine which combination
strategy generalizes best on unseen data.

\section{Dataset Description}
\label{sec:org0a89716}
\begin{longtable}{|p{6cm}|p{8cm}|}
Dataset Name & Wisconsin Diagnostic Breast Cancer (WDBC)\\
Dataset Source & UCI ML Repository / \texttt{sklearn.datasets.load\_breast\_cancer}\\
Number of Samples & 569\\
Number of Features & 30 numerical attributes\\
Number of Classes & 2 (Malignant, Benign)\\
Missing Values & None\\
Train-Test Split & 80--20, stratified\\
\end{longtable}

\section{Setup and Imports}
\label{sec:orgf1d81c1}

\begin{verbatim}
import numpy as np
import pandas as pd
import matplotlib.pyplot as plt
import seaborn as sns

from sklearn.datasets import load_breast_cancer
from sklearn.model_selection import train_test_split, GridSearchCV, learning_curve, StratifiedKFold
from sklearn.preprocessing import StandardScaler
from sklearn.tree import DecisionTreeClassifier
from sklearn.svm import SVC
from sklearn.naive_bayes import GaussianNB
from sklearn.linear_model import LogisticRegression
from sklearn.ensemble import BaggingClassifier, AdaBoostClassifier, GradientBoostingClassifier, StackingClassifier
from sklearn.metrics import (accuracy_score, precision_score, recall_score, f1_score,
			      confusion_matrix, ConfusionMatrixDisplay, roc_curve, auc,
			      RocCurveDisplay, classification_report)

sns.set_style("whitegrid")
plt.rcParams["figure.figsize"] = (7, 5)
RANDOM_STATE = 42
np.random.seed(RANDOM_STATE)
print("Ready")
\end{verbatim}

\begin{verbatim}
data = load_breast_cancer()
df = pd.DataFrame(data.data, columns=data.feature_names)
df["target"] = data.target  # 0 = malignant, 1 = benign
target_names = data.target_names
print("Dataset shape:", df.shape)
\end{verbatim}

\section{Exploratory Data Analysis}
\label{sec:org67c1ca1}
Include all relevant visualizations.
\begin{itemize}
\item Dataset overview
\item Statistical summary
\item Missing value analysis
\item Class distribution
\item Correlation matrix
\item Feature distribution
\end{itemize}

\begin{verbatim}
print(df.describe().T.to_string())
print("\nMissing values total:", df.isnull().sum().sum())
\end{verbatim}

\begin{verbatim}
Traceback (most recent call last):
  File "<stdin>", line 1, in <module>
  File "/tmp/babel-3a4bfv/python-gkAtsB", line 1, in <module>
    print(df.describe().T.to_string())
NameError: name 'df' is not defined
\end{verbatim}


\begin{verbatim}
plt.figure(figsize=(5,4))
ax = sns.countplot(x=df["target"].map({0:"Malignant", 1:"Benign"}))
plt.title("Class Distribution")
plt.xlabel("Diagnosis")
plt.ylabel("Count")
for p in ax.patches:
    ax.annotate(int(p.get_height()), (p.get_x()+p.get_width()/2, p.get_height()), ha="center", va="bottom")
plt.tight_layout()
plt.savefig("class_distribution.png")
plt.close()
"class_distribution.png"
\end{verbatim}

\begin{center}
\includegraphics[width=.9\linewidth]{class_distribution.png}
\end{center}

\textbf{Inference:} The dataset shows a moderate class imbalance, with more benign than malignant
samples. This should be kept in mind when interpreting raw accuracy, and is why precision,
recall, and F1-score are also reported alongside accuracy for every model.

\begin{verbatim}
plt.figure(figsize=(14,12))
corr = df.drop(columns=["target"]).corr()
sns.heatmap(corr, cmap="coolwarm", center=0, square=True, cbar_kws={"shrink":0.6})
plt.title("Correlation Matrix of Features")
plt.tight_layout()
plt.savefig("correlation_matrix.png")
plt.close()
"correlation_matrix.png"
\end{verbatim}

\begin{center}
\includegraphics[width=.9\linewidth]{correlation_matrix.png}
\end{center}

\textbf{Inference:} Several groups of features (e.g. mean radius, mean perimeter, mean area) are highly
correlated since they are geometric derivatives of one another. This multicollinearity is well
tolerated by tree-based and margin-based ensemble members but is useful context when interpreting
feature importance.

\begin{verbatim}
cols_to_plot = ["mean radius", "mean texture", "mean perimeter", "mean area",
		"mean smoothness", "mean concavity"]
fig, axes = plt.subplots(2, 3, figsize=(15, 8))
for ax, col in zip(axes.ravel(), cols_to_plot):
    sns.histplot(data=df, x=col, hue=df["target"].map({0:"Malignant",1:"Benign"}), kde=True, ax=ax)
    ax.set_title(col)
plt.tight_layout()
plt.savefig("feature_distributions.png")
plt.close()
"feature_distributions.png"
\end{verbatim}

\begin{center}
\includegraphics[width=.9\linewidth]{feature_distributions.png}
\end{center}

\textbf{Inference:} Malignant and benign samples show visibly different distributions for several mean
features (larger radius/perimeter/area/concavity tend to be malignant), confirming these features
carry discriminative signal for the ensembles built below.

\section{Data Preprocessing}
\label{sec:org6d7d341}
\begin{itemize}
\item Missing value handling: none required (no missing values in WDBC).
\item Encoding: target is already numeric (0/1); no categorical encoding needed.
\item Feature scaling: \texttt{StandardScaler} applied (fit on train, applied to test) since SVM and
Logistic Regression in the stacking ensemble are scale-sensitive.
\item Train-test split: 80--20, stratified on the target to preserve class balance.
\item Feature engineering: none applied; all 30 original features are used directly.
\end{itemize}

\begin{verbatim}
X = df.drop(columns=["target"])
y = df["target"]

X_train, X_test, y_train, y_test = train_test_split(
    X, y, test_size=0.20, stratify=y, random_state=RANDOM_STATE
)

scaler = StandardScaler()
X_train_scaled = scaler.fit_transform(X_train)
X_test_scaled = scaler.transform(X_test)

print("Train shape:", X_train.shape, " Test shape:", X_test.shape)
\end{verbatim}

\begin{verbatim}
Traceback (most recent call last):
  File "<stdin>", line 1, in <module>
  File "/tmp/babel-3a4bfv/python-aLwCB5", line 1, in <module>
    X = df.drop(columns=["target"])
NameError: name 'df' is not defined
\end{verbatim}

\section{Machine Learning Algorithm}
\label{sec:org4b36340}

\subsection{Bagging (Bootstrap Aggregation)}
\label{sec:org92c3c52}
\begin{itemize}
\item Working principle: trains many independent copies of a base estimator (Decision Tree here) on
different bootstrap resamples of the training data, then aggregates predictions by majority vote.
\item Mathematical formulation: for \(B\) bootstrap samples, the ensemble prediction is
\(\hat{y} = \text{mode}\{h_1(x), h_2(x), \ldots, h_B(x)\}\), where each \(h_b\) is trained on a
resample \(D_b\) drawn with replacement from \(D\).
\item Advantages: reduces variance, parallelizable, works well with high-variance base learners.
\item Limitations: does little to reduce bias; large ensembles increase inference cost.
\end{itemize}

\subsection{Boosting (AdaBoost / Gradient Boosting)}
\label{sec:org959f713}
\begin{itemize}
\item Working principle: builds an additive model sequentially, where each new weak learner focuses on
the samples the previous ensemble misclassified (AdaBoost re-weights samples; Gradient Boosting
fits new learners to the residual/gradient of the loss).
\item Mathematical formulation: \(F_m(x) = F_{m-1}(x) + \gamma_m h_m(x)\), where \(h_m\) is chosen to
minimize the loss given the current ensemble \(F_{m-1}\).
\item Advantages: reduces bias, often achieves very high accuracy.
\item Limitations: sequential (harder to parallelize), sensitive to noisy data/outliers, more prone to
overfitting if not regularized (learning rate, depth, number of estimators).
\end{itemize}

\subsection{Stacked Ensemble (Stacking)}
\label{sec:orga5cb30f}
\begin{itemize}
\item Working principle: trains several heterogeneous base learners (SVM, Naive Bayes, Decision Tree),
then trains a meta-learner (Logistic Regression) on their out-of-fold predictions to learn the
optimal way to combine them.
\item Mathematical formulation: \(\hat{y} = g(h_1(x), h_2(x), \ldots, h_K(x))\), where \(g\) is the
meta-learner and \(h_1, \ldots, h_K\) are the base learners.
\item Advantages: exploits complementary strengths of diverse models; often outperforms any single base
learner or a homogeneous ensemble.
\item Limitations: more complex and computationally expensive; risk of overfitting the meta-learner if
base-learner predictions aren't produced out-of-fold (mitigated here via internal CV in
\texttt{StackingClassifier}).
\end{itemize}

\section{Experimental Setup}
\label{sec:org2f9b72e}
\begin{center}
\begin{tabular}{ll}
Python Version & \\
Scikit-Learn Version & \\
IDE & \\
Operating System & \\
Hardware & \\
\end{tabular}
\end{center}

\begin{verbatim}
import sys, sklearn, platform
print("Python:", sys.version.split()[0])
print("scikit-learn:", sklearn.__version__)
print("Platform:", platform.platform())
\end{verbatim}

\begin{verbatim}
A module that was compiled using NumPy 1.x cannot be run in
NumPy 2.2.6 as it may crash. To support both 1.x and 2.x
versions of NumPy, modules must be compiled with NumPy 2.0.
Some module may need to rebuild instead e.g. with 'pybind11>=2.12'.

If you are a user of the module, the easiest solution will be to
downgrade to 'numpy<2' or try to upgrade the affected module.
We expect that some modules will need time to support NumPy 2.

Traceback (most recent call last):  File "<stdin>", line 1, in <module>
  File "/tmp/babel-3a4bfv/python-pvQAjd", line 1, in <module>
    import sys, sklearn, platform
  File "/home/emacs/.local/lib/python3.10/site-packages/sklearn/__init__.py", line 73, in <module>
    from .base import clone  # noqa: E402
  File "/home/emacs/.local/lib/python3.10/site-packages/sklearn/base.py", line 20, in <module>
    from .utils._missing import is_scalar_nan
  File "/home/emacs/.local/lib/python3.10/site-packages/sklearn/utils/__init__.py", line 9, in <module>
    from ._chunking import gen_batches, gen_even_slices
  File "/home/emacs/.local/lib/python3.10/site-packages/sklearn/utils/_chunking.py", line 11, in <module>
    from ._param_validation import Interval, validate_params
  File "/home/emacs/.local/lib/python3.10/site-packages/sklearn/utils/_param_validation.py", line 14, in <module>
    from scipy.sparse import csr_matrix, issparse
  File "/usr/lib/python3/dist-packages/scipy/sparse/__init__.py", line 267, in <module>
    from ._csr import *
  File "/usr/lib/python3/dist-packages/scipy/sparse/_csr.py", line 10, in <module>
    from ._sparsetools import (csr_tocsc, csr_tobsr, csr_count_blocks,
AttributeError: _ARRAY_API not found
Traceback (most recent call last):
  File "<stdin>", line 1, in <module>
  File "/tmp/babel-3a4bfv/python-pvQAjd", line 1, in <module>
    import sys, sklearn, platform
  File "/home/emacs/.local/lib/python3.10/site-packages/sklearn/__init__.py", line 73, in <module>
    from .base import clone  # noqa: E402
  File "/home/emacs/.local/lib/python3.10/site-packages/sklearn/base.py", line 20, in <module>
    from .utils._missing import is_scalar_nan
  File "/home/emacs/.local/lib/python3.10/site-packages/sklearn/utils/__init__.py", line 9, in <module>
    from ._chunking import gen_batches, gen_even_slices
  File "/home/emacs/.local/lib/python3.10/site-packages/sklearn/utils/_chunking.py", line 11, in <module>
    from ._param_validation import Interval, validate_params
  File "/home/emacs/.local/lib/python3.10/site-packages/sklearn/utils/_param_validation.py", line 14, in <module>
    from scipy.sparse import csr_matrix, issparse
  File "/usr/lib/python3/dist-packages/scipy/sparse/__init__.py", line 267, in <module>
    from ._csr import *
  File "/usr/lib/python3/dist-packages/scipy/sparse/_csr.py", line 10, in <module>
    from ._sparsetools import (csr_tocsc, csr_tobsr, csr_count_blocks,
ImportError: numpy.core.multiarray failed to import
\end{verbatim}

\section{Hyperparameter Selection}
\label{sec:orga1c9820}

\subsection{Bagging}
\label{sec:org9762444}
\begin{verbatim}
bagging_param_grid = {
    "n_estimators": [10, 25, 50, 100],
    "max_samples": [0.5, 0.7, 1.0],
    "max_features": [0.5, 0.7, 1.0],
}

bagging_base = BaggingClassifier(
    estimator=DecisionTreeClassifier(random_state=RANDOM_STATE),
    random_state=RANDOM_STATE
)

cv5 = StratifiedKFold(n_splits=5, shuffle=True, random_state=RANDOM_STATE)

bagging_grid = GridSearchCV(
    bagging_base, bagging_param_grid, cv=cv5,
    scoring={"accuracy": "accuracy", "f1": "f1"}, refit="accuracy", n_jobs=-1
)
bagging_grid.fit(X_train_scaled, y_train)

bagging_results = pd.DataFrame(bagging_grid.cv_results_)[
    ["param_n_estimators", "param_max_samples", "param_max_features",
     "mean_test_accuracy", "mean_test_f1"]
].rename(columns={
    "param_n_estimators": "n_estimators",
    "param_max_samples": "max_samples",
    "param_max_features": "max_features",
    "mean_test_accuracy": "Avg CV Accuracy",
    "mean_test_f1": "Avg CV F1 Score",
})
bagging_results["Avg CV Accuracy (%)"] = (bagging_results["Avg CV Accuracy"] * 100).round(2)
bagging_results = bagging_results[["n_estimators", "max_samples", "max_features",
				    "Avg CV Accuracy (%)", "Avg CV F1 Score"]]
print(bagging_results.sort_values("Avg CV Accuracy (%)", ascending=False).head(10).to_string(index=False))
print("\nBest Bagging params:", bagging_grid.best_params_)
print("Best Bagging CV accuracy: {:.4f}".format(bagging_grid.best_score_))
best_bagging = bagging_grid.best_estimator_
\end{verbatim}

\begin{verbatim}
Traceback (most recent call last):
  File "<stdin>", line 1, in <module>
  File "/tmp/babel-3a4bfv/python-QdybUm", line 7, in <module>
    bagging_base = BaggingClassifier(
NameError: name 'BaggingClassifier' is not defined
\end{verbatim}

\subsection{Boosting}
\label{sec:orgd3212ec}
\begin{verbatim}
ada_param_grid = {
    "n_estimators": [50, 100, 150],
    "learning_rate": [0.01, 0.1, 0.5, 1.0],
}
ada_base = AdaBoostClassifier(random_state=RANDOM_STATE)
ada_grid = GridSearchCV(ada_base, ada_param_grid, cv=cv5,
			 scoring={"accuracy": "accuracy", "f1": "f1"}, refit="accuracy", n_jobs=-1)
ada_grid.fit(X_train_scaled, y_train)
ada_results = pd.DataFrame(ada_grid.cv_results_)[
    ["param_n_estimators", "param_learning_rate", "mean_test_accuracy", "mean_test_f1"]
].rename(columns={"param_n_estimators": "n_estimators", "param_learning_rate": "learning_rate",
		   "mean_test_accuracy": "Avg CV Accuracy", "mean_test_f1": "Avg CV F1 Score"})
ada_results["Avg CV Accuracy (%)"] = (ada_results["Avg CV Accuracy"] * 100).round(2)
ada_results = ada_results[["n_estimators", "learning_rate", "Avg CV Accuracy (%)", "Avg CV F1 Score"]]
print("=== AdaBoost ===")
print(ada_results.sort_values("Avg CV Accuracy (%)", ascending=False).head(10).to_string(index=False))
print("Best params:", ada_grid.best_params_, " Best CV accuracy: {:.4f}".format(ada_grid.best_score_))
\end{verbatim}

\begin{verbatim}
Traceback (most recent call last):
  File "<stdin>", line 1, in <module>
  File "/tmp/babel-3a4bfv/python-gKPF2m", line 5, in <module>
    ada_base = AdaBoostClassifier(random_state=RANDOM_STATE)
NameError: name 'AdaBoostClassifier' is not defined
\end{verbatim}


\begin{verbatim}
gb_param_grid = {
    "n_estimators": [50, 100, 150],
    "learning_rate": [0.01, 0.1, 0.2],
    "max_depth": [2, 3, 4],
}
gb_base = GradientBoostingClassifier(random_state=RANDOM_STATE)
gb_grid = GridSearchCV(gb_base, gb_param_grid, cv=cv5,
			scoring={"accuracy": "accuracy", "f1": "f1"}, refit="accuracy", n_jobs=-1)
gb_grid.fit(X_train_scaled, y_train)
gb_results = pd.DataFrame(gb_grid.cv_results_)[
    ["param_n_estimators", "param_learning_rate", "param_max_depth", "mean_test_accuracy", "mean_test_f1"]
].rename(columns={"param_n_estimators": "n_estimators", "param_learning_rate": "learning_rate",
		   "param_max_depth": "max_depth", "mean_test_accuracy": "Avg CV Accuracy",
		   "mean_test_f1": "Avg CV F1 Score"})
gb_results["Avg CV Accuracy (%)"] = (gb_results["Avg CV Accuracy"] * 100).round(2)
gb_results = gb_results[["n_estimators", "learning_rate", "max_depth", "Avg CV Accuracy (%)", "Avg CV F1 Score"]]
print("=== Gradient Boosting ===")
print(gb_results.sort_values("Avg CV Accuracy (%)", ascending=False).head(10).to_string(index=False))
print("Best params:", gb_grid.best_params_, " Best CV accuracy: {:.4f}".format(gb_grid.best_score_))

if ada_grid.best_score_ >= gb_grid.best_score_:
    best_boosting = ada_grid.best_estimator_
    best_boosting_name = "AdaBoost"
else:
    best_boosting = gb_grid.best_estimator_
    best_boosting_name = "Gradient Boosting"
print("Selected boosting model:", best_boosting_name)
\end{verbatim}

\begin{verbatim}
Traceback (most recent call last):
  File "<stdin>", line 1, in <module>
  File "/tmp/babel-3a4bfv/python-qSMjts", line 6, in <module>
    gb_base = GradientBoostingClassifier(random_state=RANDOM_STATE)
NameError: name 'GradientBoostingClassifier' is not defined
\end{verbatim}

\subsection{Stacked Ensemble}
\label{sec:orgd8459e7}
\begin{verbatim}
stack_base = StackingClassifier(
    estimators=[
	("svm", SVC(probability=True, random_state=RANDOM_STATE)),
	("nb", GaussianNB()),
	("dt", DecisionTreeClassifier(random_state=RANDOM_STATE)),
    ],
    final_estimator=LogisticRegression(max_iter=1000, random_state=RANDOM_STATE),
    cv=5
)

stack_param_grid_full = {
    "svm__C": [0.1, 1, 10],
    "svm__kernel": ["linear", "rbf"],
    "dt__max_depth": [3, 5, None],
    "final_estimator__C": [0.1, 1, 10],
}

stack_grid = GridSearchCV(stack_base, stack_param_grid_full, cv=cv5,
			   scoring={"accuracy": "accuracy", "f1": "f1"}, refit="accuracy", n_jobs=-1)
stack_grid.fit(X_train_scaled, y_train)

stack_results = pd.DataFrame(stack_grid.cv_results_)[
    ["param_svm__C", "param_svm__kernel", "param_dt__max_depth", "param_final_estimator__C",
     "mean_test_accuracy", "mean_test_f1"]
].rename(columns={"param_svm__C": "SVM C", "param_svm__kernel": "SVM kernel",
		   "param_dt__max_depth": "DT max_depth", "param_final_estimator__C": "LogReg C",
		   "mean_test_accuracy": "Avg CV Accuracy", "mean_test_f1": "Avg CV F1 Score"})
stack_results["Avg CV Accuracy (%)"] = (stack_results["Avg CV Accuracy"] * 100).round(2)
stack_results = stack_results[["SVM C", "SVM kernel", "DT max_depth", "LogReg C",
				"Avg CV Accuracy (%)", "Avg CV F1 Score"]]
print(stack_results.sort_values("Avg CV Accuracy (%)", ascending=False).head(10).to_string(index=False))
print("\nBest Stacking params:", stack_grid.best_params_)
print("Best Stacking CV accuracy: {:.4f}".format(stack_grid.best_score_))
best_stacking = stack_grid.best_estimator_
\end{verbatim}

\begin{verbatim}
Traceback (most recent call last):
  File "<stdin>", line 1, in <module>
  File "/tmp/babel-3a4bfv/python-64rz6A", line 1, in <module>
    stack_base = StackingClassifier(
NameError: name 'StackingClassifier' is not defined
\end{verbatim}

\section{Results}
\label{sec:org0ae6d71}

\begin{center}
\begin{tabular}{lc}
Metric & Value\\
\hline
Accuracy & \\
Precision & \\
Recall & \\
F1-score & \\
ROC-AUC & \\
\end{tabular}
\end{center}

\textbf{(Fill using the best-performing model's row from the Performance Comparison table produced below.)}

\begin{verbatim}
models = {
    "Bagging": best_bagging,
    f"Boosting ({best_boosting_name})": best_boosting,
    "Stacked Ensemble": best_stacking,
}

results_table = []
roc_data = {}

for name, model in models.items():
    y_pred = model.predict(X_test_scaled)
    y_proba = model.predict_proba(X_test_scaled)[:, 1]
    acc = accuracy_score(y_test, y_pred)
    prec = precision_score(y_test, y_pred)
    rec = recall_score(y_test, y_pred)
    f1 = f1_score(y_test, y_pred)
    fpr, tpr, _ = roc_curve(y_test, y_proba)
    roc_auc = auc(fpr, tpr)
    roc_data[name] = (fpr, tpr, roc_auc)
    results_table.append({
	"Model": name, "Accuracy (%)": round(acc * 100, 2), "Precision": round(prec, 4),
	"Recall": round(rec, 4), "F1 Score": round(f1, 4), "ROC-AUC": round(roc_auc, 4),
    })

performance_df = pd.DataFrame(results_table)
print(performance_df.to_string(index=False))
\end{verbatim}

\begin{verbatim}
Traceback (most recent call last):
  File "<stdin>", line 1, in <module>
  File "/tmp/babel-3a4bfv/python-O0pJXL", line 2, in <module>
    "Bagging": best_bagging,
NameError: name 'best_bagging' is not defined
\end{verbatim}

\section{Result Visualization}
\label{sec:orgba60bbb}
Include all graphs generated during the experiment. For every graph: figure, caption, and 3--5
line inference.

Recommended figures: Correlation Matrix (above), Class Distribution (above), Confusion Matrix, ROC
Curve, Learning Curve. (Feature Importance is not directly applicable to the Naive Bayes / SVM
components of the stack and is omitted; Precision-Recall curves can be added analogously to the
ROC curve code below if required.)

\begin{verbatim}
fig, axes = plt.subplots(1, 3, figsize=(16, 4.5))
for ax, (name, model) in zip(axes, models.items()):
    y_pred = model.predict(X_test_scaled)
    cm = confusion_matrix(y_test, y_pred)
    ConfusionMatrixDisplay(cm, display_labels=target_names).plot(ax=ax, colorbar=False, cmap="Blues")
    ax.set_title(name)
plt.tight_layout()
plt.savefig("confusion_matrices.png")
plt.close()
"confusion_matrices.png"
\end{verbatim}

\begin{center}
\includegraphics[width=.9\linewidth]{confusion_matrices.png}
\end{center}

\textbf{Inference:} Compare the false-negative counts (malignant misclassified as benign) across the
three models -- this is the clinically costlier error. The model with the fewest false negatives
here should be preferred even if its overall accuracy is marginally lower than another model's.

\begin{verbatim}
plt.figure(figsize=(6.5, 6))
for name, (fpr, tpr, roc_auc) in roc_data.items():
    plt.plot(fpr, tpr, label=f"{name} (AUC = {roc_auc:.3f})")
plt.plot([0, 1], [0, 1], "k--", linewidth=1)
plt.xlabel("False Positive Rate")
plt.ylabel("True Positive Rate")
plt.title("ROC Curves -- Bagging vs Boosting vs Stacking")
plt.legend(loc="lower right")
plt.tight_layout()
plt.savefig("roc_curves.png")
plt.close()
"roc_curves.png"
\end{verbatim}

\begin{center}
\includegraphics[width=.9\linewidth]{roc_curves.png}
\end{center}

\textbf{Inference:} The model whose curve hugs the top-left corner most closely, and has the highest
AUC, separates the two classes best across all decision thresholds, not just the default 0.5
cutoff used for the confusion matrix above.

\begin{verbatim}
fig, axes = plt.subplots(1, 3, figsize=(17, 4.5))
for ax, (name, model) in zip(axes, models.items()):
    train_sizes, train_scores, val_scores = learning_curve(
	model, X_train_scaled, y_train, cv=5, scoring="accuracy",
	train_sizes=np.linspace(0.1, 1.0, 6), n_jobs=-1
    )
    ax.plot(train_sizes, train_scores.mean(axis=1), "o-", label="Training score")
    ax.plot(train_sizes, val_scores.mean(axis=1), "o-", label="Cross-validation score")
    ax.set_title(name)
    ax.set_xlabel("Training set size")
    ax.set_ylabel("Accuracy")
    ax.legend()
plt.tight_layout()
plt.savefig("learning_curves.png")
plt.close()
"learning_curves.png"
\end{verbatim}

\begin{center}
\includegraphics[width=.9\linewidth]{learning_curves.png}
\end{center}

\textbf{Inference:} A large, persistent gap between the training and cross-validation curves indicates
high variance (typical of an under-regularized Bagging ensemble); curves that converge to a
lower plateau together indicate higher bias (undertrained Boosting). The gap should shrink as
training size grows for a well-tuned model.

\begin{verbatim}
for name, model in models.items():
    print(f"\n=== {name} ===")
    print(classification_report(y_test, model.predict(X_test_scaled), target_names=target_names))
\end{verbatim}

\begin{verbatim}
Traceback (most recent call last):
  File "<stdin>", line 1, in <module>
  File "/tmp/babel-3a4bfv/python-yksYVn", line 1, in <module>
    for name, model in models.items():
NameError: name 'models' is not defined. Did you mean: 'codecs'?
\end{verbatim}

\section{Discussion}
\label{sec:orge5a9f7f}
\begin{itemize}
\item Bagging reduces \textbf{variance} by averaging many independently-trained high-variance trees over
bootstrap resamples; it does not directly target bias.
\item Boosting reduces \textbf{bias} by sequentially focusing on the hardest, previously misclassified
samples, at some risk of overfitting if left unregularized.
\item Stacking benefits from heterogeneous base learners (SVM, Naive Bayes, Decision Tree) because each
makes different kinds of errors; the meta-learner can learn to weight each base learner's
strengths, typically outperforming any single member and often the homogeneous ensembles too.
\item Compare the strengths/limitations/improvements you observed against the Performance Comparison
table and the confusion matrices/ROC curves above once you've run the notebook.
\end{itemize}

\section{Conclusion}
\label{sec:org8e03cf4}
Bagging, Boosting, and Stacked Ensemble models were implemented on the WDBC dataset and tuned via
5-fold cross-validation. Based on the test-set metrics in the Performance Comparison table, state
which ensemble generalized best and why (typically Stacking, since it combines complementary
inductive biases from SVM, Naive Bayes, and Decision Tree base learners under a Logistic Regression
meta-learner).

\section{References}
\label{sec:org853bc7a}
\begin{itemize}
\item Scikit-learn: Ensemble Methods, \url{https://scikit-learn.org/stable/modules/ensemble.html}
\item Scikit-learn: \texttt{BaggingClassifier}, \url{https://scikit-learn.org/stable/modules/generated/sklearn.ensemble.BaggingClassifier.html}
\item UCI Machine Learning Repository: Breast Cancer Wisconsin (Diagnostic) Data Set, \url{https://archive.ics.uci.edu/dataset/17}
\end{itemize}
\end{document}
